\documentclass[a4paper,12pt]{article}

\usepackage[utf8]{inputenc}
\usepackage[russian]{babel}
\usepackage[left=2cm,right=2cm,top=2cm,bottom=2.5cm]{geometry}

\usepackage{amsmath}
\usepackage{amssymb}
\usepackage{booktabs}
\usepackage{tabularx}
\usepackage{caption}
\usepackage{enumitem}
\usepackage{hyperref}
\usepackage{cmap}

\newcolumntype{Y}{>{\raggedleft\arraybackslash}X}

\DeclareMathOperator{\EL}{EL}
\DeclareMathOperator{\PD}{PD}
\DeclareMathOperator{\DD}{DD}
\DeclareMathOperator{\CVaR}{CVaR}
\DeclareMathOperator{\RAROC}{RAROC}
\DeclareMathOperator{\CRC}{CRC}

\setlength{\parindent}{0pt}
\setlength{\parskip}{6pt plus 2pt minus 1pt}
\renewcommand{\baselinestretch}{1.15}

\hypersetup{
    colorlinks=true,
    linkcolor=black,
    urlcolor=blue,
    citecolor=black,
}

\title{\vspace{-1.5cm}\textbf{Модель управления кредитным портфелем\\
    с учётом влияния макроэкономических факторов}}
% \author{\textbf{Дмитрий КУРЕННОЙ}, Т-Банк \\[2pt]
%     \textbf{Максим КИРЯКИН}, ВМК МГУ имени М.~В.~Ломоносова}
\author{}
\date{}

\begin{document}

\maketitle
\thispagestyle{empty}

\textbf{Классические системы мониторинга кредитного риска опираются на финансовую отчётность
заёмщиков, которая публикуется с существенным запаздыванием и не позволяет банку реагировать
на ухудшение кредитного качества портфеля до того, как риск реализуется. В статье на примере
портфеля из четырнадцати публичных российских компаний показано, как объединить структурную
модель оценки вероятности дефолта, сценарное прогнозирование макрофакторов и многокритериальную
оптимизацию структуры портфеля в единую систему непрерывного мониторинга. Приведены результаты
12-месячного walk-forward-бэктеста четырёх стратегий управления.}

\bigskip

\textbf{Ключевые слова:} кредитный риск, дистанция до дефолта, модель Мертона, стресс-тестирование,
макроэкономические сценарии, портфельная оптимизация, CVaR, Risk Parity, RAROC.

\bigskip

\section{Ограничения статического контроля кредитного риска}

На кредитный риск приходится около 70\,\% совокупного риск-профиля российского банковского
сектора, поэтому именно качество его мониторинга определяет устойчивость кредитной организации.
Тем не менее в большинстве банков система раннего предупреждения по-прежнему построена вокруг
периодического пересмотра финансовых коэффициентов заёмщика: долговой нагрузки, ликвидности,
рентабельности. Эти показатели описывают прошлое состояние компании и обновляются не чаще
одного раза в квартал. Между двумя датами отчётности заёмщик для модели остаётся неизменным,
хотя его реальное кредитное качество меняется ежедневно вместе с рыночной конъюнктурой.

Санкционные ограничения, кризисные эпизоды 2020 и 2022 годов, резкие колебания ключевой ставки
и курса рубля наглядно показали недостаточность статического подхода: повышение ставки и
ослабление рубля меняют кредитное качество заёмщиков задолго до того, как это попадёт
в очередную отчётность. Возникает практическая потребность в системе мониторинга, которая
переоценивает риск портфеля непрерывно и связывает кредитное качество каждого заёмщика
с прогнозом макросреды.

Ниже описан подход, объединяющий три блока в единую систему: структурную оценку вероятности
дефолта по рыночным данным, эконометрическую связь кредитного качества с макрофакторами и
оптимизационную ребалансировку портфеля. Именно их интеграция в замкнутый цикл принятия
решений представляет практическую ценность для банка.

\section{Структурная оценка вероятности дефолта в реальном времени}

Первый блок отвечает на вопрос, насколько близок заёмщик к дефолту в текущий момент. Для публичных
компаний на него позволяет ответить структурная модель Мертона, лежащая в основе методологии
Moody's KMV. Дефолт в ней трактуется как ситуация, при которой рыночная стоимость активов
компании $V$ опускается ниже порога обязательств $D$. Ключевая метрика~--- дистанция до дефолта
$\DD$, измеряющая, на сколько стандартных отклонений стоимость активов превышает долговой порог:

\[
    \DD = \frac{\ln\!\left(\dfrac{V}{D}\right) + \left(r - \dfrac{\sigma_A^2}{2}\right) T}
              {\sigma_A \sqrt{T}},
\]

где $r$~--- безрисковая ставка, $\sigma_A$~--- волатильность активов, $T$~--- горизонт анализа.
Перевод дистанции до дефолта в вероятность дефолта выполняется через функцию стандартного
нормального распределения:

\[
    \PD = \Phi(-\DD).
\]

Главное практическое достоинство этого блока для мониторинга~--- частота обновления. Стоимость и
волатильность активов восстанавливаются из ежедневных биржевых котировок, поэтому $\DD$ и $\PD$
пересчитываются с той же частотой, не дожидаясь публикации отчётности. Ожидаемые потери по позиции считаются стандартно:

\[
    \EL_{i} = \PD_{i} \cdot \mathrm{LGD}_{i} \cdot \mathrm{EAD}_{i},
\]

где $\mathrm{LGD}_i$~--- доля потерь при дефолте, $\mathrm{EAD}_i$~--- сумма под риском.

У подхода есть ограничения: он применим к публичным компаниям с ликвидными акциями, а
модель Мертона не учитывает, что в фазе рецессии дефолты коррелируют между собой.
Поэтому одной структурной модели недостаточно~--- нужен учёт макроэкономической составляющей.

\section{Связь кредитного качества с макросредой}

Второй блок связывает дистанцию до дефолта с состоянием экономики. Для этого на панельных данных
строится регрессия дистанции до дефолта на набор макрофакторов~--- инфляцию, ключевую ставку,
уровень безработицы и курс USD/RUB:

\[
    \DD_{i,t} = \alpha_i + \beta_1 \, \text{Инфляция}_t + \beta_2 \, \text{Ставка}_t
              + \beta_3 \, \text{Безработица}_t + \beta_4 \, \text{Курс}_t + \varepsilon_{i,t},
\]

На эмпирической базе из публичных компаний пяти секторов экономики (нефтегазового, финансового,
металлургического, телекоммуникационного и потребительского) совместная значимость инфляции,
ключевой ставки и уровня безработицы подтверждена F-тестом с $p < 10^{-16}$. Знаки коэффициентов
экономически интерпретируемы: рост инфляции, ключевой ставки и безработицы снижает дистанцию до
дефолта, то есть увеличивает кредитный риск. Статистическая значимость этих связей позволяет
банку опираться на макропрогноз при оценке кредитного качества портфеля на формальной
количественной основе, а не только на экспертных суждениях.

\section{Сценарное прогнозирование макрофакторов}

Чтобы оценивать риск с опережением, недостаточно знать текущую связь
кредитного качества с макросредой~--- нужен прогноз самой макросреды. Третий блок строит сценарные
траектории макрофакторов тремя взаимодополняющими моделями временных рядов:

\begin{itemize}[leftmargin=2em, itemsep=2pt]
    \item \textbf{VAR}~--- векторная авторегрессия, учитывающая взаимные связи макропоказателей;
    \item \textbf{SARIMAX}~--- модель с сезонностью и внешними регрессорами;
    \item \textbf{Prophet}~--- аддитивная модель тренда и сезонности.
\end{itemize}

На walk-forward-бэктесте по метрикам MAE, RMSE и MAPE наименьшие ошибки практически по всем
макрорядам показала SARIMAX, поэтому она выбрана базовой моделью прогноза. Прогнозы VAR и Prophet
используются как альтернативные сценарии при стресс-тестировании. Если на исторических данных
модели дают близкие оценки, то заметный разброс их прогнозов на будущее служит для
риск-мониторинга индикатором возросшей макроэкономической неопределённости.

Прогноз макрофакторов подставляется в регрессию из предыдущего раздела, давая прогнозную
дистанцию до дефолта и вероятность дефолта каждого заёмщика на горизонте нескольких
месяцев вперёд~--- то есть оценку риска с опережением, а не по факту его реализации.

\section{Оптимизация структуры портфеля и стратегии управления}

Знание прогнозной вероятности дефолта позволяет не только наблюдать риск, но и управлять им.
Четвёртый блок~--- многокритериальная оптимизация структуры портфеля. В практике кредитного
комитета задача ставится не в долях капитала, а в абсолютных суммах: банк располагает бюджетом
$B$ и рассматривает $N$ заявок, по каждой из которых принимает решение об одобряемой сумме
$x_i \in [0, a_i]$. Каждая заявка $i$ описывается запрашиваемой суммой $a_i$, ставкой $r_i^{cr}$,
индивидуальным $\mathrm{LGD}_i$ и сроком $T_i$. Годовая вероятность дефолта из модели Мертона
приводится к кумулятивной за срок кредита, $\PD_i^{cum} = 1 - (1 - \PD_i^{ann})^{T_i/12}$,
что корректно отражает накопление риска во времени.

Целевая функция объединяет три компоненты~--- риск-скорректированную доходность, меру рыночного
риска и ожидаемые кредитные потери~--- при бюджетном, индивидуальных и секторальных ограничениях:

\[
    J(x) = -\lambda_R \, \frac{\sum_{i} x_i R_i}{B}
           + \lambda_V \, \rho(w)
           + \lambda_{EL} \sum_{i} w_i \, \PD_i^{cum} \, \mathrm{LGD}_i \, \frac{12}{T_i}
           \;\to\; \min_{x},
\]
\[
    \sum_{i} x_i \le B, \qquad 0 \le x_i \le a_i, \qquad
    \sum_{i \in S_j} x_i \le c_j B,
\]

где $w_i = x_i / \sum_j x_j$~--- портфельные веса, $R_i = r_i^{cr}(1 - \PD_i^{cum})
- \PD_i^{cum}\,\mathrm{LGD}_i\,\tfrac{12}{T_i}$~--- риск-скорректированная доходность на рубль
вложений в заёмщика $i$, $S_j$~--- множество заёмщиков сектора $j$, $c_j$~--- лимит на сектор. Коэффициенты
$\lambda_R$, $\lambda_V$, $\lambda_{EL}$ задают баланс между доходностью, рыночным и кредитным
риском. Зависимость весов от сумм делает задачу нелинейной, поэтому она решается методом SLSQP
для каждой из трёх мер риска $\rho(w)$:

\begin{itemize}[leftmargin=2em, itemsep=2pt]
    \item \textbf{Mean-EL}~--- волатильность портфеля $\rho_{\text{vol}}(w) = \sqrt{w^\top \Sigma_V w}$
        по ковариационной матрице рыночной стоимости активов;
    \item \textbf{CVaR}~--- эмпирическая условная стоимость под риском $\CVaR_\alpha$ по сценариям
        скорректированных доходностей, контролирующая потери в неблагоприятном хвосте распределения;
    \item \textbf{Risk Parity}~--- выравнивание совокупных рисковых вкладов заёмщиков.
\end{itemize}

Стратегия Risk Parity обобщена на кредитный риск через совокупный рисковый вклад
$\CRC_i = \tfrac{w_i (\Sigma_V w)_i}{\sigma_p} + w_i \, \EL_i^{ann}$, объединяющий рыночную и кредитную
компоненты, что не допускает скрытой концентрации кредитного риска на одном заёмщике.

По результатам оптимизации рассчитывается показатель $\RAROC = \sum_i x_i^\ast R_i \big/
\sum_i x_i^\ast$, характеризующий доходность портфеля на единицу размещённого капитала с учётом
кредитных потерь. В такой постановке мониторинг помогает принимать решения: он показывает,
как изменятся профиль портфеля и его доходность при одобрении конкретной заявки.

\section{Результаты бэктеста}

Все четыре стратегии (включая пассивный бенчмарк с равными долями) сопоставлены на 12-месячном
walk-forward-бэктесте в режиме работы с абсолютными суммами заявок. Параметры эксперимента:
бюджет 20\,млрд\,руб., 14 заявок, секторальный лимит 30\,\%, порог ребалансировки $\delta = 0{,}02$,
ставка транзакционных издержек $c = 0{,}001$, коэффициенты целевой функции $\lambda_R = 0{,}2$,
$\lambda_V = 0{,}7$, $\lambda_{EL} = 0{,}1$ (приоритет отдан контролю риска). Сводные результаты
приведены в таблице~\ref{tab:strategies}.

\begin{table}[ht]
    \centering
    \captionsetup{justification=raggedright, singlelinecheck=false}
    \caption{Сравнение стратегий управления на 12-месячном бэктесте}
    \label{tab:strategies}
    \small
    \begin{tabularx}{\textwidth}{lYYYY}
        \toprule
        \textbf{Стратегия} & Нетто-доход,\,\% & $\EL$,\,\% & Волатильность,\,\% & Sharpe \\
        \midrule
        Passive (Equal, B\&H) & $-14{,}01$ & $0{,}29$ & $14{,}26$ & $0{,}21$ \\
        Mean-EL               & $-12{,}19$ & $0{,}24$ & $13{,}82$ & $0{,}26$ \\
        CVaR                  & $-11{,}48$ & $0{,}26$ & $13{,}45$ & $0{,}29$ \\
        \textbf{Risk Parity}  & $\mathbf{-10{,}94}$ & $\mathbf{0{,}22}$ & $\mathbf{12{,}98}$ & $\mathbf{0{,}33}$ \\
        \bottomrule
    \end{tabularx}
\end{table}

В рассмотренный интервал российский фондовый рынок снижался, поэтому нетто-доход всех стратегий
отрицателен. Однако активные стратегии демонстрируют заметно меньшую просадку за счёт
перераспределения капитала в пользу заёмщиков с лучшим кредитным профилем и меньшей реализованной
волатильностью. По коэффициенту Шарпа ранжирование устойчиво: $0{,}26$, $0{,}29$ и $0{,}33$ против
$0{,}21$ у пассивного портфеля. Лучший результат показывает Risk Parity: минимальная волатильность
$12{,}98\,\%$, наименьшие ожидаемые потери $0{,}22\,\%$ и наибольший Sharpe при умеренном числе
ребалансировок.

Для практики важны два наблюдения. Границы весов ($w_{\min} = 0{,}02$,
$w_{\max} = 0{,}25$) устранили вырожденные решения с концентрацией капитала, а пороговая
ребалансировка ограничила число пересмотров портфеля четырьмя-шестью за год.

\section{Практические выводы для риск-мониторинга}

Описанный подход показывает, что банк может перейти от периодического к непрерывному
мониторингу кредитного риска и оценивать его с опережением, используя только открытые
рыночные данные и общедоступные инструменты анализа. Из работы следует несколько
практических выводов для риск-подразделения банка.

\begin{enumerate}[leftmargin=2em, itemsep=3pt]
    \item \textbf{Рыночные данные дополняют отчётность.} Структурная модель Мертона переоценивает
        кредитное качество публичных заёмщиков ежедневно, закрывая разрыв между датами
        публикации отчётности.
    \item \textbf{Макропрогноз должен быть встроен в мониторинг.} Связь дистанции до дефолта
        с инфляцией, ставкой и безработицей позволяет переводить макросценарии в прогноз
        кредитного риска портфеля.
    \item \textbf{Выбор стратегии зависит от риск-аппетита.} Единая постановка с тремя
        мерами риска (волатильность, CVaR, Risk Parity) позволяет банку выбирать стратегию
        под свой профиль.
    \item \textbf{Пороговая ребалансировка экономит издержки.} Пересмотр портфеля только при
        значимом изменении прогнозного риска удерживает транзакционные издержки под контролем.
\end{enumerate}

Предложенный подход не заменяет действующие модели PD и LGD, а дополняет их как
инструмент динамического портфельного мониторинга для стресс-тестирования кредитных
портфелей, поддержки решений кредитного комитета и оценки риск-скорректированной доходности.
Полученные выводы относятся к одному 12-месячному окну, и их обобщение требует проверки
на нескольких временных интервалах и в разных состояниях рынка. Однако основной вывод
остаётся в силе: явный учёт прогнозного кредитного риска в структуре портфеля снижает
просадку и волатильность по сравнению с пассивным управлением.

\end{document}
